\documentclass[11pt,spanish,a4paper]{article}
% Versión 2.o cuat 2015 Víctor Bettachini < bettachini@df.uba.ar >

\usepackage{babel}
\addto\shorthandsspanish{\spanishdeactivate{~<>}}
\usepackage[utf8]{inputenc}
\usepackage{float}

\usepackage{units}
\usepackage[separate-uncertainty=true, multi-part-units=single, locale=FR]{siunitx}

\usepackage{amsmath}
\usepackage{amstext}
\usepackage{amssymb}

\newcommand{\pvec}[1]{\vec{#1}\mkern2mu\vphantom{#1}}

% \usepackage{tikz}
% \input{DimLinesTikz}
% \usetikzlibrary{decorations.pathmorphing, patterns}

\usepackage{graphicx}
\graphicspath{{./graphs/}}
\usepackage{wrapfig}


\usepackage[margin=1.3cm,nohead]{geometry}
% \voffset-3.5cm
% \hoffset-3cm
% \setlength{\textwidth}{17.5cm}
% \setlength{\textheight}{27cm}

\usepackage{lastpage}
\usepackage{fancyhdr}
\pagestyle{fancyplain}
\fancyhead{}
\fancyfoot{{\tiny \textcopyright Departamento de Física, FCEyN, UBA}}
\fancyfoot[C]{ {\tiny Actualizado al \today} }
\fancyfoot[RO, LE]{Pág. \thepage/\pageref{LastPage}}
\renewcommand{\headrulewidth}{0pt}
\renewcommand{\footrulewidth}{0pt}

% lista multicolumna
\usepackage{multicol}

% textsubscript
\usepackage{changes}


% \def \materia {Física II para químicos}
\def \periodo {cuatrimestre de verano - 2017}
\def \website {http://materias.df.uba.ar/f2qa2017v}


\begin{document}
\noindent
\textbf{Física II (Químicos)}\hfill \textcopyright {\tt DF, FCEyN, UBA}
% \textbf{\materia}\hfill \periodo
\begin{center}
  \textsc{\large Difracción}
  % \textsc{\large Guía 9: Difracción}
\par\end{center}{\large \par}


\begin{enumerate}

\section*{Rendija única}

\item
\begin{minipage}[t]{0.55\textwidth}
El sistema óptico de la figura está compuesto por dos lentes convergentes y una apertura rectangular muy larga de ancho \(a\).
Dé una expresión para la distribución de luz en el plano focal de \(L_2\) y, teniéndola en cuenta, diga dónde está ubicado el pico del máximo principal de difracción en los casos:
\begin{enumerate}
	\item fuente en \(f_1\),
	\item fuente en P.
\end{enumerate}
\end{minipage}
\begin{minipage}[c][1em][t]{0.35\textwidth}
    \includegraphics[width=\textwidth]{g09e01}
\end{minipage}

Halle la ubicación de la imagen geométrica de una fuente en \(f_1\) y de otra en \(P\).
¿Qué relación tiene la ubicación de dichas imágenes con la de los máximos hallados en a)?



%%%%% Hasta acá


\item
Considere la figura de difracción de Fraunhofer producida por una rendija de ancho \(a\) y largo \(b\) (\(b \gg  a\)) ubicada entre dos lentes convergentes y centrada en el eje óptico del sistema.
La fuente puntual monocromática de longitud de onda \(\lambda\) se coloca en el foco objeto de la primera lente.
\begin{enumerate}
	\item ¿Dónde se coloca la pantalla de observación para observar difracción de Fraunhofer?
	\item Calcule la posición de los máximos y de los mínimos de intensidad, el ancho angular de la campana principal de difracción y de los máximos secundarios.
	\item Calcule la relación de intensidades entre el máximo principal y el primer máximo secundario.
	\item Grafique la intensidad sobre la pantalla. ¿En función de qué variables lo hace? ¿Podría haber elegido otras? ¿Cuáles?
	\item Discuta cómo se modifican los parámetros de la figura de difracción si se cambia:
	\begin{enumerate}
		\item el ancho de la ranura.
		\item la longitud de onda.
		\item la fuente monocromática por una policromática.
	\end{enumerate}
	\item Resuelva todo el problema nuevamente si la fuente se encuentra en el plano focal objeto de la primera lente a una altura h del eje óptico.
\end{enumerate}


\item
Una rendija de ancho \(a = \SI{0.25}{\milli\metre}\) y largo \(b \gg a\) está colocada delante de una lente convergente.
richa rendija está iluminada por ondas planas que inciden sobre ella, siendo \(\lambda= \SI{500}{\nano\metre}\).
En el plano focal imagen de la lente se observa una figura de difracción.
La distancia entre el primer mínimo a la izquierda del máximo principal y el tercer mínimo a su derecha es \SI{3}{\milli\metre}.
Además, el primer mínimo a la izquierda está ubicado \SI{3}{\milli\metre} a la derecha del eje óptico.
\begin{enumerate}
	\item ¿Cuánto vale la distancia focal de la lente usada?
	\item ¿Dónde el máximo principal? ¿Dónde se encuentra la fuente?
\end{enumerate}


\section*{Doble rendija}


\item
Se tienen dos rendijas iguales de ancho \(a\), cuya separación entre centros es \(d\), colocadas entre dos lentes delgadas convergentes, ubicadas en forma simétrica respecto del eje óptico del sistema.
Una fuente puntual monocromática se encuentra en el foco de la primera lente.
Considere la figura de interferencia-difracción de Fraunhofer de la fuente.
\begin{enumerate}
	\item Calcule la posición de los máximos y mínimos tanto de interferencia como de difracción.
	\item Grafique la intensidad sobre la pantalla. ¿En función de qué variable lo hace? ¿Qué otra variable podría haber usado?
	\item ¿Cuántos órdenes de interferencia hay dentro de la campana principal de difracción?
	\item ¿Por qué motivo cuando se estudia el experimento de Young de interferencia no se tiene en cuenta el efecto de difracción en cada ranura?
\end{enumerate}


\item
Se realiza una experiencia de difracción por doble rendija con una fuente que emite en \SI{400}{\nano\metre}.
La separación entre los puntos medios de las rendijas es de \SI{0,4}{\milli\metre} y el ancho de cada una de ellas es de \SI{0,04}{\milli\metre}.
La pantalla está a \SI{1}{\metre} de las rendijas.
Si se cambia la fuente por otra que emite en \SI{600}{\nano\metre}, determine:
\begin{enumerate}
	\item en cuánto varió la interfranja,
	\item en cuánto varió el número total de franjas de interferencia contenidas en el máximo principal de difracción,
	\item y en cuánto varió el ancho angular de la campana principal de difracción.
\end{enumerate}


\item
Sobre dos ranuras de Young separadas una distancia de \SI{1}{\milli\metre} incide la superposición de dos ondas planas monocromáticas de longitudes de onda \(\lambda_1\) y \(\lambda_2\).
\begin{enumerate}
	\item ¿Qué relación debe satisfacer el cociente \(\frac{\lambda_1}{\lambda_2}\) para que el tercer orden de interferencia constructiva de \(\lambda_1\) coincida con el tercer mínimo de \(\lambda_2\)? 
	\item ¿Qué ancho deben tener las ranuras para que además esos órdenes coincidan con el primer mínimo de difracción de \(\lambda_1\)?
	\item ¿Qué intensidad se registrará en la pantalla en ese punto?
\end{enumerate}



\section*{N rendijas}


\item 
% \item \textbf{N rendijas de ancho conocido}
Una onda plana monocromática de longitud de onda \(\lambda\) incide normalmente sobre una red de transmisión plana formada por \(N\) rendijas de ancho \(a\) y de período \(d\) (\(a \ll d\)).
Suponga que la teoría escalar corresponde a una descripción exacta del fenómeno.
\begin{enumerate}
\item Se coloca una pantalla de observación de modo de observar difracción de Fraunhofer.
	Analice la distribución de intensidad sobre la pantalla y grafíquela.
\item Calcule la posición angular de las líneas espectrales y su intensidad.
	¿Las líneas espectrales corresponden a los máximos de interferencia o de difracción?
	Calcule la separación angular entre dos líneas consecutivas.
\item ¿Cuántos mínimos de interferencia hay entre dos líneas espectrales consecutivas?
	¿Cuántos máximos secundarios hay entre dichas líneas?
\item Calcule el ancho angular de las líneas espectrales.
\item Si la incidencia sobre la red no es normal, ¿cómo cambia la figura de interferencia-difracción?
\end{enumerate}



\item 
Se incide con un haz de luz monocromático de longitud de onda \(\lambda= \SI{589}{\nano\metre}\) de forma tal que se observa el siguiente patrón de difracción sobre una pantalla ubicada a \SI{150}{\centi\metre} de una red de rendijas.
Determine:
\begin{enumerate}
	\item el número de rendijas que tiene la red,
	\item la distancia entre rendijas,
	\item y el ancho de cada rendija.
\end{enumerate}
\begin{center}
	\includegraphics[width=0.75\textwidth]{g09e11}
\end{center}


\section*{Redes de difracción}

\item 
% \item \textbf{N fuentes puntuales}
Se tienen \(N\) fuentes puntuales monocromáticas de longitud de onda \(\lambda\) alineadas y separadas entre sí una distancia \(d\).
Se coloca una pantalla a una distancia \(L\) (\(L \gg d\)).
\begin{enumerate}
	\item Determine la intensidad luminosa sobre la pantalla en función de la coordenada \(y\).
	\item Determine la separación entre los máximos principales.
	\item ¿Cuántos máximos secundarios aparecen entre los máximos principales? ¿Cuántos mínimos?
	\item ¿Qué sucede al variar la cantidad de ranuras o la separación de las mismas?
	\item Si las fuentes emiten en dos longitudes de onda (\(\lambda\) y \(\lambda'\)), ¿en qué condiciones quedan nítidamente separados sus respectivos máximos?
\end{enumerate}


\item
Se dispone de dos redes de difracción por transmisión cuadradas de \SI{2}{\centi\metre} de lado, una tiene \num{600} líneas/\si{\milli\metre} y la otra 1200 líneas/\si{\milli\metre}.
Calcule:
\begin{enumerate}
	\item el poder resolvente de cada red en el primer orden,
	\item el máximo orden observable si la fuente emite en \SI{500}{\nano\metre}.
	\item ¿Es importante tener en cuenta el ángulo de incidencia?
	Considere primero que la luz incide normalmente a la red y luego repita el cálculo para el caso en que la incidencia es rasante.
	\item Calcule el máximo poder resolvente de cada una.
	\item Determine si alguna de ellas resuelve entre las longitudes de onda \(\lambda_1= \SI{500}{\nano\metre}\) y \(\lambda_2= \SI{500,007}{\nano\metre}\).
\end{enumerate}


\item
Luz de longitudes de onda \(\lambda_1= \SI{500}{\nano\metre}\) y \(\lambda_2= \SI{520}{\nano\metre}\) incide normalmente sobre una red de difracción de rendijas.
La distancia entre rendijas es de \SI{1}{\micro\metre} y se emplea una lente convergente de distancia focal de \SI{2}{\metre} para focalizar el espectro sobre una pantalla.
Calcule aproximadamente la separación lineal sobre la pantalla entre las líneas espectrales de primer orden de \(\lambda_1\) y \(\lambda_2\).




\end{enumerate}
\end{document}
